\chapter{Studiu bibliografic}\label{ch:studiubib}

\pagestyle{fancy}

Pentru realizarea unui instrument local de adnotare, documentarea bibliografică a fost organizată pe topicuri relevante domeniului, rezumate și corelate cu obiectivele proiectului.

În raportul tehnic din 2017 al lui Breitenfeld și Muller-Birn se prezintă o trecere în revistă a instrumentelor de adnotare semantică folosite în contexte de cercetare, cu accent pe legarea conținutului la vocabulari controlate și baze de cunoștințe~\cite{Breitenfeld2017}. Autorii compară uneltele după publicul țintă, formatele suportate, sursele de cunoștințe (de tip RDF/SKOS), modul de interacțiune și contextul de rulare (web, desktop sau integrare în infrastructuri de cercetare). Analiza evidențiază diferențe între unelte orientate spre utilizatori tehnici și cele care încearcă să includă experți de domeniu prin interfețe mai prietenoase.

Concluzia principală este că multe instrumente sunt orientate către utilizatori tehnici, iar suportul pentru colaborare și provenance rămâne limitat, ceea ce îngreunează adopția în echipe interdisciplinare. Deși lucrarea tratează adnotarea semantică (nu vizuală), observațiile despre uzabilitate, colaborare și trasabilitate sunt relevante pentru proiectul de față, unde aceleași nevoi apar în procesul de etichetare a imaginilor. Din această perspectivă, studiul oferă un context util pentru alegerea unui flux de lucru local, simplificat și orientat pe utilizator.

Pentru sincronizarea colaborativă, Sun et al. analizează diferențele reale dintre OT (Operational Transformation) și CRDT (Commutative Replicated Data Type) într-un cadru unificat de transformare~\cite{Sun2020OTCRDT}. Autorii separă starea vizibilă de utilizator (external state) de structurile interne de sincronizare, comparând modul în care fiecare paradigmă gestionează operațiile concurente. Ei arată că și CRDT implică transformări indirecte și costuri suplimentare, ceea ce poate afecta performanța în sisteme interactive. Concluziile sugerează că, pentru aplicații cu server central și colaborare în timp real, un model operațional bazat pe evenimente poate fi mai predictibil și mai eficient decât un CRDT complex. Pentru un instrument de adnotare colaborativă, această perspectivă susține propagarea operațiilor (de tip create/update/delete) în locul trimiterii stării complete.

Lucrarea despre Microsoft COCO introduce un set de date de referință pentru detecție, segmentare și descriere de imagini în contexte cotidiene, punând accent pe scene complexe cu obiecte multiple~\cite{lin2015microsoftcococommonobjects}. Dataset-ul include sute de mii de imagini și peste două milioane de instanțe etichetate, acoperind zeci de clase obiect, grupate în supercategorii. Anotările sunt detaliate (instance segmentation, bounding box-uri și capturi textuale), iar evaluarea folosește metrici bazate pe precizie la mai multe praguri IoU (Intersection over Union), pentru a reflecta mai realist calitatea localizării. COCO a devenit un reper pentru comparații între modele și a standardizat modul de evaluare în sarcini de adnotare și segmentare.

Pentru versionarea și trasabilitatea adnotărilor, standardul W3C PROV-DM oferă un model formal pentru a descrie cum sunt generate și modificate entitățile digitale~\cite{Moreau2013PROVDM}. Modelul introduce concepte precum entitate, activitate și agent, precum și relații de tipul wasGeneratedBy și wasDerivedFrom, care permit reconstruirea istoricului unei adnotări. În contextul anotării colaborative, aceste mecanisme susțin auditarea modificărilor, reproducibilitatea și controlul calității, mai ales atunci când adnotările sunt create atât de utilizatori, cât și de modele automate.

Calitatea adnotărilor este un factor critic pentru performanța modelelor, iar sistemele moderne folosesc mecanisme de asigurare a calității precum acordul între adnotatori, etichetarea prin consens, fluxuri de review și audituri periodice. Aceste practici reduc erorile umane și cresc consistența, mai ales în domenii sensibile, unde deciziile automate pot avea impact ridicat. Un alt aspect important este bias-ul seturilor de date, care poate distorsiona evaluările și generalizarea în practică.

În lucrarea „Unbiased Look at Dataset Bias”, Torralba și Efros analizează bias-ul seturilor de date din recunoașterea obiectelor și arată că modelele învață adesea semnături specifice dataset-ului, nu concepte vizuale generale~\cite{5995347}. Ei compară mai multe colecții (PASCAL VOC, ImageNet, LabelMe, SUN09, Caltech101, MSRC) și folosesc experimente de tip „Name That Dataset” și teste de generalizare între dataset-uri. Rezultatele indică faptul că acuratețea scade semnificativ atunci când antrenarea și testarea se fac pe surse diferite, iar bias-ul de selecție, captură și negative set afectează direct robustețea modelelor.

Studiul este pre-deep learning și se bazează pe HOG/BoW (Histogram of Oriented Gradients/Bag of Words) cu SVM (Support Vector Machine), ceea ce limitează comparațiile cu modele moderne, dar concluziile despre bias rămân relevante. Pentru un instrument de adnotare colaborativă, lucrarea susține importanța fluxurilor de QA (Quality Assurance), a consensului între adnotatori și a controlului asupra distribuției datelor și a exemplelor negative.

În raportul tehnic despre învățarea activă, Settles descrie cum un model poate obține performanțe ridicate cu mai puține etichete dacă selectează instanțele cele mai informative pentru adnotare~\cite{settles.tr09}. Lucrarea acoperă scenarii precum sinteza interogărilor, eșantionarea din flux și eșantionarea dintr-un rezervor, precum și strategii de interogare (incertitudine, comitet, reducerea erorii așteptate). În practică, învățarea activă reduce efortul de etichetare, dar introduce provocări legate de costuri variabile, oracole zgomotoase și bias indus de selecție.

Pentru un instrument colaborativ de adnotare, aceste idei sunt utile deoarece permit prioritizarea imaginilor sau a regiunilor care aduc maxim de informație, evitând munca redundantă. Abordări precum batch-mode active learning pot distribui eficient sarcinile către mai mulți adnotatori, iar mecanismele de reetichetare pot compensa erorile umane. Totuși, costul computațional al selecției și dependența de modelul curent rămân limitări care trebuie luate în calcul.

În fluxurile human-in-the-loop pentru segmentare, modelul propune contururi, iar utilizatorul corectează și validează rezultatele, reducând timpul de adnotare fără a pierde controlul asupra calității. Aceste sisteme sunt relevante în domenii cu cerințe ridicate de precizie, unde verificarea umană rămâne esențială.

Polygon-RNN propune adnotarea instanțelor ca poligoane generate secvențial de un RNN (Recurrent Neural Network) peste caracteristici extrase de un CNN (Convolutional Neural Network), iar utilizatorul poate corecta vârfurile în timpul generării~\cite{castrejon2017annotatingobjectinstancespolygonrnn}. Modelul învață să prezică vârfuri pe o grilă discretizată și se oprește la un token de final, ceea ce transformă desenarea completă într-un proces de ajustare. Autorii raportează accelerări semnificative față de adnotarea manuală, menținând un acord ridicat în termeni de IoU.

Pentru proiect, această abordare susține ideea de asistență IA (Inteligenta Artificială) cu feedback uman, unde interfața trebuie să permită corecturi rapide și să integreze imediat rezultatul în forma finală. Metoda evidențiază că un workflow mixt poate reduce efortul, dar depinde de un design UI (User Interface) clar și de modele suficient de stabile pentru a nu genera multe corecții.

ScribbleBox oferă un exemplu direct de constrângeri de timp real în adnotarea video, raportând explicit latențele modulelor sale~\cite{chen2020scribbleboxinteractiveannotationframework}. Propagarea măștilor la 512px ajunge la ~295 ms/cadru (perceptibil), iar la 256px coboară la ~152 ms; cu cache de caracteristici, timpul scade la ~82 ms, apropiindu-se de interacțiune fluentă. În schimb, propagarea scribble-urilor este rapidă (15--28 ms), iar operațiile de curve fitting și interacțiune box sunt sub 15 ms, ceea ce permite feedback aproape instant după corecții locale.

Implicațiile pentru un instrument colaborativ sunt clare: corecțiile locale trebuie să fie sincronizate (sub 100 ms), iar propagarea globală poate rula asincron, cu indicator de progres și posibilitatea de a continua adnotarea altor obiecte. Un cache persistent al trăsăturilor (feature cache) și al stărilor intermediare este esențial, iar o strategie „rapid/precis” (preview la 256px, rafinare la 512px) poate ascunde latențele. În medii multi-utilizator, aceste timpi impun cozi și prioritizare, iar pe hardware limitat devin necesare modele mai compacte sau offloading către cloud.

Panoptic segmentation unifică segmentarea semantică și cea pe instanțe într-o singură sarcină, astfel încât fiecare pixel primește o categorie semantică și, unde este cazul, un identificator de instanță. Kirillov et al. propun această formulare pentru a obține o înțelegere coerentă a scenei și introduc un criteriu de evaluare dedicat, Panoptic Quality (PQ), care combină calitatea segmentării și a recunoașterii~\cite{kirillov2019panopticsegmentation}. Lucrarea a standardizat evaluarea prin adnotări panoptice și a influențat arhitecturi ulterioare (ex. Panoptic FPN, UPSNet, modele bazate pe transformere), fiind relevantă pentru proiect prin faptul că evidențiază necesitatea unor adnotări consistente și complete la nivel de imagine.

Instance segmentation combină detecția obiectelor cu segmentarea semantică, dar diferențiază instanțele distincte ale aceleiași clase. În lucrarea „Simultaneous Detection and Segmentation”, Hariharan et al. propun un cadru unificat care prezice simultan bounding box-uri și măști pentru fiecare obiect, folosind reprezentări pe regiuni și clasificare la nivel de instanță~\cite{hariharan2014simultaneousdetectionsegmentation}. Metoda extinde detectoarele bazate pe regiuni prin integrarea unei predicții de mască, oferind o înțelegere mai fină a scenei decât segmentarea semantică clasică.

Relevanța pentru proiect este directă, deoarece adnotarea de instanțe cere atât localizarea obiectelor, cât și delimitarea precisă a contururilor. Lucrarea oferă un precedent conceptual pentru fluxuri de lucru care combină desenarea de bounding box-uri cu rafinarea măștilor, ceea ce susține nevoia de instrumente de adnotare flexibile și eficiente.

U-Net abordează segmentarea semantică în contexte cu puține date adnotate, propunând o arhitectură complet convoluțională în formă de U, cu cale de contractare și expansiune și conexiuni de tip skip pentru păstrarea detaliilor de rezoluție înaltă~\cite{ronneberger2015unetconvolutionalnetworksbiomedical}. Modelul folosește o strategie overlap-tile pentru imagini mari, augmentare cu deformări elastice și o funcție de pierdere ponderată care accentuează separarea obiectelor adiacente. În evaluările ISBI, rezultatele au depășit metodele bazate pe ferestre glisante, raportând scoruri IoU ridicate chiar cu seturi de antrenare mici.

Limitările includ faptul că ieșirea este mai mică decât intrarea (convoluții fără umplutură), iar antrenarea necesită atenție la inițializare și la dimensiunea lotului. Pentru proiect, U-Net susține ideea de auto-segmentare rapidă în interfață, antrenare eficientă cu puține exemple și procesare pe dale pentru imagini de rezoluții mari.

CLIP (Contrastive Language-Image Pre-training) propune un mod de învățare a conceptelor vizuale din limbaj natural, antrenând simultan un encoder vizual și unul textual pe un set foarte mare de perechi imagine-text~\cite{radford2021learningtransferablevisualmodels}. Obiectivul contrastiv maximizează similaritatea pentru perechi corecte și o minimizează pentru cele greșite, ceea ce permite transfer zero-shot: etichetele țintă sunt transformate în prompturi textuale (de tip „A photo of a \{label\}.”), iar imaginea este clasată prin potrivire cu descrierea cea mai probabilă. Rezultatele arată că CLIP poate concura cu modele complet supervizate pe ImageNet fără antrenare specifică pe acel dataset și este mai robust la schimbări de distribuție (schițe, randări) decât modelele antrenate strict pe ImageNet. Encoderul vizual CLIP a devenit un fundament pentru modele open-vocabulary care adaptează reprezentările la regiuni mascate în segmentare~\cite{radford2021learningtransferablevisualmodels,liang2023openvocabularysemanticsegmentationmaskadapted}.

Lucrarea evidențiază și rolul prompt engineering: formulările contextuale și ensembling-ul de prompturi cresc performanța zero-shot. Limitările includ dificultăți în clasificare fine-grained, numărare și sarcini abstracte, precum și dependența de un volum masiv de date; în plus, bias-urile din datele web se pot reflecta în predicții, cu implicații etice și de confidențialitate~\cite{radford2021learningtransferablevisualmodels}. Pentru proiect, CLIP este relevant ca fundament pentru interogări textuale flexibile în pre-adnotare și pentru că poate oferi sugestii rapide fără antrenare suplimentară, dar impune nevoia de verificare umană și de control al bias-urilor în fluxul colaborativ.

DistilBERT propune distilarea cunoștințelor direct în faza de pre-antrenare, folosind un model profesor BERT și un student mai mic care imită distribuțiile de ieșire ale profesorului (soft labels) și alinierea embedding-urilor, printr-o funcție de pierdere compusă (distillation + language modeling + cosine)~\cite{sanh2020distilbertdistilledversionbert}. Rezultatul este un model cu 40\% mai puțini parametri, cu inferență mai rapidă și performanță comparabilă (păstrează majoritatea capacității de înțelegere). În contextul chatbot-ului din platformă, DistilBERT este relevant pentru rularea locală sau pe edge a sarcinilor de NLP (completări, clasificări de intenții, sugestii de etichete), reducând latența și costurile, păstrând totodată posibilitatea de a apela modele mai mari doar când este nevoie.

Lucrarea despre segmentarea semantică open-vocabulary cu Mask-Adapted CLIP pornește de la limitarea seturilor închise și de la faptul că CLIP scade mult pe imagini mascate, din cauza domain gap-ului și a „zero tokens” introduși de fundal~\cite{liang2023openvocabularysemanticsegmentationmaskadapted}. Autorii adaptează CLIP pe regiuni mascate și descrieri textuale, construind un set de date zgomotos din COCO Captions (extragere de substantive + aliniere automată a măștilor cu MaskFormer și CLIP) și propun Mask Prompt Tuning (MPT), care înlocuiește direct tokenii de fundal mascat cu tokeni de prompt învățabili, fie cu CLIP înghețat, fie cu fine-tuning complet.

OVSeg obține SOTA (State-of-the-Art) zero-shot (ex. mIoU 29.6 pe ADE20K-150, +8.5 față de OpenSeg, recorduri pe ADE20K-847 și Pascal Context-459), iar MPT singur aduce un câștig consistent cu puțini parametri antrenabili~\cite{liang2023openvocabularysemanticsegmentationmaskadapted}. Limitările includ propuneri de măști influențate de clasele văzute, ambiguități lingvistice în evaluare, decalaj între clase văzute/nevăzute și faptul că modelele supervizate moderne rămân superioare. Pentru proiect, tehnica susține interogări textuale libere, pre-adnotare automată din metadate și adaptare eficientă în medii colaborative, cu un cost redus de memorie/GPU și nevoie de review pentru ambiguități lexicale.

YOLOE (modelul propus; extensie a familiei YOLO, You Only Look Once) oferă un model unificat, open-set (clase deschise) și real-time pentru detecție și segmentare, integrând prompturi textuale, vizuale și un mod prompt-free în aceeași arhitectură~\cite{wang2025yoloerealtimeseeing}. Pentru prompturi textuale, folosește RepRTA (Re-parameterizable Region-Text Alignment, aliniere regiune-text re-parametrizabilă): un bloc ușor de tip FFN (Feed-Forward Network) rafinează embedding-urile textuale la antrenare, apoi este re-parametrizat în capul de clasificare ca o convoluție $1\times1$, fără cost suplimentar la inferență. Pentru prompturi vizuale, SAVPE (Semantic-Activated Visual Prompt Encoder, encoder de prompturi vizuale activat semantic) separă o ramură semantică (caracteristici agnostice de prompt) de o ramură de activare (ponderi specifice promptului), agregându-le eficient. În mod prompt-free, LRPC (Lazy Region-Prompt Contrast, contrast regiune-prompt „leneș”) învață un embedding care localizează toate obiectele, iar asocierea cu un vocabular extins se face doar pe ancorele relevante.

Rezultatele arată performanță open-set competitivă cu eficiență ridicată: reducere de ~3x a costului de antrenare față de YOLO-World, câștiguri pe LVIS (Large Vocabulary Instance Segmentation) în zero-shot (fără adaptare pe setul de test) și viteză mai mare la inferență, plus transfer solid pe COCO față de YOLOv8~\cite{wang2025yoloerealtimeseeing}. Limitările includ dependența de un vocabular static în modul prompt-free, de un encoder textual extern și necesitatea generării de pseudo-măști cu SAM (Segment Anything Model) pentru antrenarea segmentării. Pentru proiect, YOLOE este relevant prin suportul multimodal al adnotării (text sau indicii vizuale), posibilitatea unui mod de auto-etichetare prompt-free eficient și faptul că re-parametrizarea elimină costuri la rulare, facilitând deployment pe servere ieftine sau pe client (edge) în fluxuri colaborative.

Mask2Former (modelul propus; Masked-attention Mask Transformer) abordează fragmentarea dintre segmentarea semantică, panoptică și de instanțe, propunând o arhitectură universală bazată pe mask classification, unde fiecare segment este reprezentat ca un object query~\cite{cheng2022maskedattentionmasktransformeruniversal}. Inovația este masked attention: în locul cross-attention global, interogările pot „vedea” doar regiunea prezisă la etapa anterioară, ceea ce accelerează convergența și reduce zgomotul de fundal. Decodorul folosește o strategie multi-scale eficientă (runde succesive pe $1/32$, $1/16$, $1/8$), rulează masked attention înainte de self-attention și înlocuiește query-urile cu parametri antrenabili, astfel încât modelul propune intern regiuni validate direct de funcția de pierdere.

Pentru memorie, pierderile sunt calculate pe un eșantion de puncte informative, nu pe toți pixelii, reducând consumul VRAM la antrenare. Rezultatele raportează SOTA simultan pe COCO panoptic (PQ), COCO instance (AP) și ADE20K semantic (mIoU), cu convergență în ~50 epoci față de sute de epoci la MaskFormer/DETR~\cite{cheng2022maskedattentionmasktransformeruniversal}. Limitările includ nevoia de antrenare separată per sarcină/dataset și performanță mai slabă pe obiecte foarte mici. Pentru proiect, Mask2Former este relevant prin unificarea infrastructurii de modele, posibilitatea de fine-tuning eficient pe seturi create colaborativ și generarea de contururi propuse cu focalizare locală, utilă pentru unelte de pre-select și corecție rapidă.

Neurosurgeon propune o schemă de „collaborative intelligence” între cloud și edge, unde un DNN este împărțit la un punct optim de separare între dispozitivul mobil și serverul cloud~\cite{10.1145/3037697.3037698}. Modelul estimează timpul de execuție și energia pe fiecare strat atât local, cât și în cloud, apoi alege stratul $L$ la care se transmite activarea către server, optimizând latența sau consumul energetic în funcție de lățimea de bandă și capabilitățile dispozitivului. Evaluările pe mai multe rețele arată că această împărțire reduce semnificativ latența sau energia față de rularea completă pe cloud sau complet locală.

Relevanța pentru proiect este directă: un instrument colaborativ de adnotare poate rula local straturile inițiale pentru feedback rapid (ex. propuneri brute), iar straturile costisitoare pot fi delegate către backend pentru rafinare. Astfel, se obține o experiență interactivă cu latență mică, în timp ce resursele serverului sunt folosite eficient, iar sistemul rămâne adaptiv la condițiile de rețea și la puterea dispozitivului utilizatorului~\cite{10.1145/3037697.3037698}.

